\subsection{Negative bundle metrics and logarithmic weights}
\label{subsec:negative-bundle-logarithmic-weights}
\index{Vector bundle!negative curvature and logarithmic weights}

For positive curvature, the reciprocal norm in
\cref{lem:reciprocal-bundle-norm-levi} is useful near a removed
zero section. With negative curvature, the squared norm itself
is plurisubharmonic. Adding its logarithm gives a weight whose
Levi form controls both the base and the fiber directions.

In this subsection the base $X$ need not be compact. Suppose
that it has a K\"ahler form $\omega$ and a smooth nonnegative
plurisubharmonic exhaustion $\rho$. Here an exhaustion has
compact sublevel sets; a constant function qualifies on a compact
base. Let $\pi\colon E\to X$ be a holomorphic vector bundle of rank
$r\geq1$, with a smooth Hermitian metric $h$, and let $Z$ be its
zero section. Use the unscaled curvature and coefficient forms
$\alpha_{h,v}$ of \cref{subsec:projective-complement-curvature-index},
and assume
\[
 \alpha_{h,v}\leq-|v|_h^2\omega
 \qquad(x\in X,\ v\in E_x).
\]
This is a uniform negative curvature bound with respect to the
chosen base form.

\begin{proposition}\label{prop:negative-bundle-logarithmic-weight}
\index{Plurisubharmonic function!negative bundle norm}
Put $U(v)=|v|_h^2$, $q=U\mathrm e^U$ and
$\Psi=\rho\circ\pi+U$. Then $\Psi$ is a smooth
plurisubharmonic exhaustion of $E$ and
\[
 \Omega=\pi^*\omega+2\pi\ddc U
\]
is a K\"ahler form on $E$. The function $\log q$, extended by
$-\infty$ on $Z$, is plurisubharmonic and satisfies
\[
 2\pi\ddc\log q\geq\Omega\quad\text{on }E\setminus Z.
\]
Both $U$ and $q$ are strictly plurisubharmonic off $Z$.
For the real radial vector field $\mathcal R$, whose flow is
$v\mapsto\mathrm e^s v$, one has
$\mathcal Rq=2(1+U)q$. In particular $q$ has no critical
points off $Z$.
\end{proposition}
\begin{proof}
At $x$, choose a Chern-normal unitary holomorphic frame as in
\cref{lem:reciprocal-bundle-norm-levi}. Write a tangent vector at
$(x,v)$ as $(a,b)$, with $a$ horizontal and $b$ vertical.
The first base derivatives and mixed base--fiber derivatives of
$U$ vanish there. Let $\mathcal L$ denote the unscaled complex
Hessian, corresponding to $2\pi\ddc=\mathrm i\partial\bar\partial$.
The normal-frame calculation of that lemma gives
\[
 \begin{aligned}
 \mathcal L_U(a,b)
   &=-h(R_h(a,\bar a)v,v)+|b|_h^2
      \geq U\omega(a,\bar a)+|b|_h^2,\\
 \mathcal L_{\log U}(a,b)
   &=-\frac{h(R_h(a,\bar a)v,v)}U
       +\frac{|b|_h^2U-|h(b,v)|^2}{U^2}
      \geq\omega(a,\bar a)\quad(v\ne0).
 \end{aligned}
\]
The last vertical term is nonnegative by Cauchy--Schwarz.
These formulas show that $U$ is plurisubharmonic, that
$\Omega$ is positive definite even on $Z$, and that
$2\pi\ddc(\log U+U)\geq\Omega$ off $Z$.
The form $\Omega$ is closed because $\omega$ is closed.
The first formula gives strict positivity for $U$ off $Z$;
exponentiating $\log q$ gives it for $q$.
Locally, $U$ is comparable to the squared Euclidean fiber norm,
so $\log q$ tends to $-\infty$ at $Z$ and is locally bounded above.
By \cref{lem:qpsh-removable-analytic}, the upper-limit extension
is plurisubharmonic and has value $-\infty$ on $Z$.

The function $\Psi$ is plurisubharmonic. Its sublevel set of
height $M$ lies over the compact set $\{\rho\leq M\}$ and in
the disk bundle $\{U\leq M\}$ over that set, which is compact.
The sublevel set is closed, so $\Psi$ is an exhaustion. Finally,
$\mathcal RU=2U$ gives the radial identity and excludes critical
points where $q>0$.
\end{proof}

\begin{corollary}\label{cor:negative-bundle-weight-jets}
\index{Multiplier ideal!zero section logarithmic weight}
For every $m>0$,
\[
 \mathcal I(m\log q)
    =\mathcal I_Z^{\max\{\lfloor m\rfloor-r+1,0\}}.
\]
In particular, if $k\geq1$ is an integer and $m\geq r+k-1$,
a holomorphic function that is locally square integrable against
$q^{-m}$ times a smooth positive volume form vanishes to order
at least $k$ along $Z$. If $m<r+k-1$, this conclusion fails for
some local holomorphic functions.
\end{corollary}
\begin{proof}
In a local bundle frame, $\log q=\log|v|^2+O(1)$ near $Z$.
Apply \cref{lem:smooth-center-multiplier-ideal} in codimension
$r$. Its radial integrability condition for a homogeneous fiber
term of degree $a$ is $a+r>m$. For $m<r+k-1$, a monomial
of degree $k-1$ satisfies this condition and has insufficient
vanishing order. Smooth positive volume forms do not change
local integrability.
\end{proof}

\begin{corollary}\label{cor:negative-bundle-complete-punctured-metric}
\index{K\"ahler metric!complete punctured negative bundle}
The manifold $E\setminus Z$ admits a complete K\"ahler metric.
\end{corollary}
\begin{proof}
Choose a smooth convex increasing function $\eta$ on $\mathbb R$
equal to $-\log(-t)$ for $t\leq-1$. On $E\setminus Z$ put
\[
 \widehat\Omega
   =\Omega+2\pi\ddc(\Psi^2)
       +2\pi\ddc\eta(\log q).
\]
The added forms are nonnegative, so this is K\"ahler.
Since $\Psi\geq0$ is plurisubharmonic,
$2\pi\ddc(\Psi^2)\geq
2\mathrm i\partial\Psi\wedge\bar\partial\Psi$.
Thus $\Psi$ has bounded gradient for $\widehat\Omega$, and a
curve of finite length stays in a compact sublevel of $\Psi$ in
$E$. The only possible escape from that compact set within
$E\setminus Z$ is approach to $Z$.
Where $\log q\leq-1$, set $\tau=\log(-\log q)$. Convexity and
the chain rule give
\[
 2\pi\ddc\eta(\log q)
   \geq\mathrm i\partial\tau\wedge\bar\partial\tau.
\]
Hence $\tau$ also has bounded gradient there, while
$\tau\to+\infty$ on approach to $Z$. A finite-length curve
cannot escape in this way either. The criterion that every
escaping curve has infinite length proves completeness.
\end{proof}

For example, let $B\to X$ have a smooth Hermitian metric with
positive unscaled curvature form $\omega$. The direct sum
$E=(B^*)^{\oplus r}$ with its induced metric satisfies
$\alpha_{h,v}=-|v|_h^2\omega$.
Since $\omega$ is closed, the construction applies whenever $X$
also has a plurisubharmonic exhaustion as above, in particular
when $X$ is compact.

If the base is connected, compact and has positive dimension,
one may take $\rho=0$. Then $U$ is an exhaustion that is
strictly plurisubharmonic off $Z$, but the total space $E$ is
not Stein: every holomorphic function on $E$ restricts to a
constant on the compact connected zero section and therefore
fails to separate its distinct points. Strict positivity away
from a compact analytic set does not remove this obstruction.
